\documentclass[12pt,a4paper]{article}

\usepackage[utf8]{inputenc}
\usepackage[T1]{fontenc}
\usepackage{amsmath, amssymb}
\usepackage{enumitem}
\usepackage{geometry}
\usepackage{xcolor}
\usepackage{multicol}

\geometry{margin=2.3cm}

\title{Banco de Questões Objetivas -- Aula 11\\Técnicas de Programação SQL}
\author{Banco de Dados}
\date{}

\begin{document}

\maketitle

\section*{Instruções}

Este banco contém questões objetivas sobre os principais tópicos da Aula 11:
\begin{itemize}
    \item programação para banco de dados;
    \item linguagem hospedeira e sublinguagem de dados;
    \item descompasso de impedância;
    \item SQL embutido;
    \item cursores;
    \item SQL dinâmica;
    \item SQLJ;
    \item JDBC;
    \item Hibernate;
    \item SQL/CLI;
    \item procedimentos e funções armazenadas;
    \item SQL/PSM;
    \item SQL Injection;
    \item dados criptografados.
\end{itemize}

\section{Questões de Múltipla Escolha}

\begin{enumerate}[label=\textbf{Q\arabic*.}]

\item Em aplicações reais, a maioria das interações com bancos de dados ocorre por meio de:

\begin{enumerate}[label=\alph*)]
    \item programas de aplicação cuidadosamente projetados e testados.
    \item digitação manual de comandos pelo administrador.
    \item arquivos de texto sem conexão com SGBD.
    \item consultas feitas apenas por planilhas.
\end{enumerate}

\item A interface interativa de um SGBD é conveniente principalmente para:

\begin{enumerate}[label=\alph*)]
    \item criação de esquema, criação de restrições e consultas ad-hoc.
    \item processamento automático de todas as aplicações web.
    \item substituição completa de programas de aplicação.
    \item eliminação de conexões com o servidor.
\end{enumerate}

\item Em programação de banco de dados, a linguagem hospedeira é:

\begin{enumerate}[label=\alph*)]
    \item a linguagem de programação geral usada pela aplicação, como Java, C ou C++.
    \item sempre a linguagem SQL.
    \item o nome físico da tabela no disco.
    \item uma restrição de integridade.
\end{enumerate}

\item Em programação de banco de dados, a sublinguagem de dados geralmente é:

\begin{enumerate}[label=\alph*)]
    \item SQL.
    \item HTML.
    \item CSS.
    \item Assembly.
\end{enumerate}

\item Uma das abordagens para programação de BD é:

\begin{enumerate}[label=\alph*)]
    \item comandos de BD embutidos em uma linguagem de programação.
    \item ausência total de linguagem hospedeira.
    \item uso exclusivo de arquivos PDF.
    \item proibição de bibliotecas de acesso.
\end{enumerate}

\item A abordagem baseada em biblioteca de funções de BD consiste em:

\begin{enumerate}[label=\alph*)]
    \item disponibilizar uma API para chamadas ao banco a partir da linguagem hospedeira.
    \item colocar SQL dentro de imagens.
    \item impedir conexão com vários bancos.
    \item substituir todas as tabelas por arquivos texto.
\end{enumerate}

\item Uma linguagem completamente nova para BD é ideal principalmente quando:

\begin{enumerate}[label=\alph*)]
    \item a aplicação possui interação intensiva com o banco de dados.
    \item não existe banco de dados.
    \item deseja-se evitar qualquer acesso ao SGBD.
    \item o banco deve ser usado apenas manualmente.
\end{enumerate}

\item O descompasso de impedância ocorre por causa da:

\begin{enumerate}[label=\alph*)]
    \item incompatibilidade entre o modelo da linguagem hospedeira e o modelo do banco.
    \item igualdade perfeita entre Java e SQL.
    \item inexistência de tipos de dados no banco.
    \item impossibilidade de usar cursores.
\end{enumerate}

\item Um exemplo de descompasso de impedância é:

\begin{enumerate}[label=\alph*)]
    \item diferenças entre tipos de dados da linguagem hospedeira e do BD.
    \item uma tabela possuir chave primária.
    \item uma consulta ter \texttt{WHERE}.
    \item o banco possuir apenas uma tabela.
\end{enumerate}

\item Resultados de consultas SQL geralmente são:

\begin{enumerate}[label=\alph*)]
    \item bags de tuplas, isto é, sequências de valores de atributos.
    \item sempre um único número inteiro.
    \item arquivos executáveis.
    \item classes Java obrigatoriamente.
\end{enumerate}

\item Para processar múltiplas tuplas retornadas por uma consulta em uma linguagem hospedeira, normalmente usa-se:

\begin{enumerate}[label=\alph*)]
    \item cursor ou iterador.
    \item comando \texttt{DROP}.
    \item chave estrangeira.
    \item criptografia MD5 obrigatória.
\end{enumerate}

\item A interação básica de um programa cliente com o BD segue a ordem:

\begin{enumerate}[label=\alph*)]
    \item abrir conexão, submeter comandos, fechar conexão.
    \item fechar conexão, submeter comandos, abrir conexão.
    \item criar índice, apagar tabela, abrir conexão.
    \item executar \texttt{DROP}, executar \texttt{DELETE}, encerrar o SGBD.
\end{enumerate}

\item SQL embutido significa que:

\begin{enumerate}[label=\alph*)]
    \item comandos SQL são inseridos dentro de uma linguagem de programação geral.
    \item comandos SQL são proibidos dentro de programas.
    \item o banco é acessado apenas por interface gráfica.
    \item o SGBD deixa de usar SQL.
\end{enumerate}

\item Em SQL embutido, os comandos SQL podem incluir:

\begin{enumerate}[label=\alph*)]
    \item definições de dados, restrições, consultas, atualizações e visões.
    \item apenas comandos \texttt{SELECT}.
    \item apenas comandos \texttt{DROP}.
    \item somente comentários.
\end{enumerate}

\item Em SQL embutido, comandos como \texttt{EXEC SQL} servem para:

\begin{enumerate}[label=\alph*)]
    \item distinguir comandos SQL dos comandos da linguagem hospedeira.
    \item apagar o banco automaticamente.
    \item declarar uma chave primária em Java.
    \item criar uma conexão HTTP.
\end{enumerate}

\item Variáveis compartilhadas em SQL embutido geralmente aparecem prefixadas por:

\begin{enumerate}[label=\alph*)]
    \item dois pontos, como \texttt{:variavel}.
    \item arroba, como \texttt{@variavel}.
    \item sustenido, como \texttt{\#variavel}.
    \item cifrão, como \texttt{\$variavel}.
\end{enumerate}

\item O pré-processador em SQL embutido pode:

\begin{enumerate}[label=\alph*)]
    \item separar comandos SQL embutidos do programa na linguagem hospedeira.
    \item remover todas as tabelas do banco.
    \item substituir o SGBD por um arquivo local.
    \item impedir a compilação da aplicação.
\end{enumerate}

\item No SQL embutido em C, variáveis declaradas em uma seção \texttt{DECLARE} são usadas para:

\begin{enumerate}[label=\alph*)]
    \item compartilhar valores entre o programa e comandos SQL.
    \item criar automaticamente todas as tabelas.
    \item impedir comunicação com o banco.
    \item apagar exceções do compilador.
\end{enumerate}

\item Em SQL embutido, variáveis como \texttt{SQLCODE} ou \texttt{SQLSTATE} são usadas para:

\begin{enumerate}[label=\alph*)]
    \item comunicar erros ou exceções entre o BD e o programa.
    \item armazenar a senha criptografada do usuário.
    \item substituir todos os cursores.
    \item criar uma nova linguagem de programação.
\end{enumerate}

\item O comando \texttt{CONNECT TO server-name AS connection-name AUTHORIZATION user-account-info;} é usado para:

\begin{enumerate}[label=\alph*)]
    \item conectar o programa a um servidor de banco de dados.
    \item fechar todas as conexões.
    \item criar uma tabela.
    \item executar uma consulta aninhada.
\end{enumerate}

\item O comando \texttt{SET CONNECTION connection-name;} é usado para:

\begin{enumerate}[label=\alph*)]
    \item mudar a conexão ativa.
    \item remover uma conexão do banco.
    \item apagar a conexão fisicamente.
    \item criar uma função armazenada.
\end{enumerate}

\item O comando \texttt{DISCONNECT connection-name;} é usado para:

\begin{enumerate}[label=\alph*)]
    \item desconectar uma conexão.
    \item criar uma conexão.
    \item mudar a conexão ativa.
    \item criar uma trigger.
\end{enumerate}

\item Em SQL embutido, quando uma consulta retorna múltiplas tuplas, é necessário usar:

\begin{enumerate}[label=\alph*)]
    \item cursor.
    \item \texttt{DROP TABLE}.
    \item chave primária composta.
    \item \texttt{DELETE} sem \texttt{WHERE}.
\end{enumerate}

\item O comando \texttt{FETCH} em um cursor serve para:

\begin{enumerate}[label=\alph*)]
    \item mover o cursor para a próxima tupla.
    \item apagar todas as tuplas.
    \item criar uma nova conexão.
    \item declarar uma procedure.
\end{enumerate}

\item O comando \texttt{CLOSE CURSOR} indica que:

\begin{enumerate}[label=\alph*)]
    \item o processamento da consulta foi completado.
    \item uma tabela foi removida.
    \item uma conexão foi aberta.
    \item uma senha foi criptografada.
\end{enumerate}

\item A SQL dinâmica permite:

\begin{enumerate}[label=\alph*)]
    \item compor e executar declarações SQL em tempo de execução.
    \item executar apenas consultas previamente compiladas.
    \item impedir qualquer comando digitado pelo usuário.
    \item substituir SQL por XML.
\end{enumerate}

\item Na SQL embutida tradicional, cada nova consulta pode exigir:

\begin{enumerate}[label=\alph*)]
    \item alteração do código-fonte.
    \item remoção do banco.
    \item desligamento do servidor.
    \item criação de uma nova linguagem.
\end{enumerate}

\item Uma atualização dinâmica costuma ser mais simples que uma consulta dinâmica porque:

\begin{enumerate}[label=\alph*)]
    \item em consultas, tipos e números de atributos recuperados podem ser desconhecidos em tempo de compilação.
    \item atualizações nunca usam SQL.
    \item consultas dinâmicas não retornam dados.
    \item atualizações sempre removem tabelas.
\end{enumerate}

\item SQLJ é:

\begin{enumerate}[label=\alph*)]
    \item um padrão para embutir SQL em Java.
    \item uma linguagem para apagar bancos.
    \item uma forma de criptografia de senha.
    \item um tipo de chave estrangeira.
\end{enumerate}

\item O tradutor SQLJ transforma comandos SQLJ em Java usando a interface:

\begin{enumerate}[label=\alph*)]
    \item JDBC.
    \item HTML.
    \item CSS.
    \item XML puro.
\end{enumerate}

\item Em SQLJ, para processar várias tuplas, são usados:

\begin{enumerate}[label=\alph*)]
    \item iteradores.
    \item comandos \texttt{DROP}.
    \item chaves estrangeiras.
    \item triggers obrigatórias.
\end{enumerate}

\item Em SQLJ, um iterador designado é associado ao resultado de uma consulta por meio de:

\begin{enumerate}[label=\alph*)]
    \item lista de tipos e nomes dos atributos.
    \item apenas nomes de tabelas.
    \item apenas senhas de usuário.
    \item apenas comandos \texttt{DELETE}.
\end{enumerate}

\item Em SQLJ, um iterador posicional lista:

\begin{enumerate}[label=\alph*)]
    \item somente os tipos dos atributos.
    \item somente os nomes dos atributos.
    \item somente as tabelas do banco.
    \item somente os usuários conectados.
\end{enumerate}

\item JDBC significa:

\begin{enumerate}[label=\alph*)]
    \item Java Database Connectivity.
    \item Java Data Backup Control.
    \item Join Database Command.
    \item Java Domain Constraint.
\end{enumerate}

\item JDBC é uma:

\begin{enumerate}[label=\alph*)]
    \item biblioteca de funções de BD disponível para Java.
    \item linguagem completamente nova que substitui Java.
    \item técnica exclusiva para apagar bancos.
    \item criptografia de senha.
\end{enumerate}

\item Um programa Java com JDBC pode acessar qualquer SGBD relacional que possua:

\begin{enumerate}[label=\alph*)]
    \item driver JDBC.
    \item arquivo PDF.
    \item comando \texttt{DROP}.
    \item senha em texto puro.
\end{enumerate}

\item JDBC permite que um programa:

\begin{enumerate}[label=\alph*)]
    \item conecte-se a vários bancos ou fontes de dados.
    \item acesse apenas um banco fixo sem driver.
    \item elimine a necessidade de SQL.
    \item funcione sem conexão.
\end{enumerate}

\item Um dos primeiros passos no acesso a BD com JDBC é:

\begin{enumerate}[label=\alph*)]
    \item importar a biblioteca \texttt{java.sql.*}.
    \item executar \texttt{DROP TABLE}.
    \item fechar a conexão antes de abrir.
    \item criptografar todas as tabelas.
\end{enumerate}

\item Em JDBC, um objeto \texttt{Connection} representa:

\begin{enumerate}[label=\alph*)]
    \item uma conexão com o banco de dados.
    \item uma tupla retornada.
    \item uma coluna da tabela.
    \item uma função armazenada.
\end{enumerate}

\item Em JDBC, um objeto \texttt{Statement} representa:

\begin{enumerate}[label=\alph*)]
    \item um comando SQL a ser executado.
    \item uma conexão física com a internet.
    \item uma tabela resultante.
    \item uma senha criptografada.
\end{enumerate}

\item Em JDBC, parâmetros de comandos SQL podem ser indicados por:

\begin{enumerate}[label=\alph*)]
    \item ponto de interrogação, \texttt{?}.
    \item asterisco, \texttt{*}.
    \item dois pontos, \texttt{:}.
    \item barra, \texttt{/}.
\end{enumerate}

\item Em JDBC, \texttt{PreparedStatement} é recomendado principalmente quando:

\begin{enumerate}[label=\alph*)]
    \item a consulta será executada múltiplas vezes ou possui parâmetros.
    \item a consulta nunca será executada.
    \item o objetivo é apagar a conexão.
    \item o banco não suporta SQL.
\end{enumerate}

\item Em JDBC, antes de executar um \texttt{PreparedStatement}, os parâmetros devem ser:

\begin{enumerate}[label=\alph*)]
    \item ligados a variáveis do programa.
    \item apagados do banco.
    \item transformados em tabelas.
    \item convertidos obrigatoriamente para XML.
\end{enumerate}

\item Em JDBC, \texttt{executeQuery} é usado principalmente para:

\begin{enumerate}[label=\alph*)]
    \item comandos de consulta que retornam \texttt{ResultSet}.
    \item comandos \texttt{INSERT}, \texttt{DELETE} e \texttt{UPDATE}.
    \item criação de classes Java.
    \item fechamento de conexão.
\end{enumerate}

\item Em JDBC, \texttt{executeUpdate} é usado principalmente para:

\begin{enumerate}[label=\alph*)]
    \item comandos \texttt{INSERT}, \texttt{DELETE} ou \texttt{UPDATE}.
    \item consultas que retornam \texttt{ResultSet}.
    \item abrir conexão.
    \item criar driver.
\end{enumerate}

\item Em JDBC, \texttt{ResultSet} pode ser entendido como:

\begin{enumerate}[label=\alph*)]
    \item uma tabela bidimensional com o resultado da consulta.
    \item uma chave primária.
    \item uma conexão ativa.
    \item uma procedure armazenada.
\end{enumerate}

\item Em JDBC, o método \texttt{next()} de um \texttt{ResultSet} serve para:

\begin{enumerate}[label=\alph*)]
    \item avançar para a próxima linha do resultado.
    \item apagar a próxima linha da tabela.
    \item fechar a conexão.
    \item criar uma nova tabela.
\end{enumerate}

\item Hibernate facilita principalmente:

\begin{enumerate}[label=\alph*)]
    \item o mapeamento objeto-relacional.
    \item a remoção de todos os bancos relacionais.
    \item a substituição de classes Java por SQL puro.
    \item a criação de senhas MD5.
\end{enumerate}

\item Hibernate pode transformar:

\begin{enumerate}[label=\alph*)]
    \item classes Java em tabelas de dados e vice-versa.
    \item consultas SQL em imagens.
    \item conexões em arquivos PDF.
    \item senhas em chaves primárias.
\end{enumerate}

\item Uma vantagem do Hibernate apresentada na aula é:

\begin{enumerate}[label=\alph*)]
    \item gerar chamadas SQL e liberar o desenvolvedor de conversões manuais.
    \item impedir portabilidade entre SGBDs.
    \item eliminar completamente a necessidade de banco de dados.
    \item substituir o SGBD por HTML.
\end{enumerate}

\item Em Hibernate, um arquivo XML pode ser usado para:

\begin{enumerate}[label=\alph*)]
    \item relacionar propriedades do objeto aos campos da tabela.
    \item apagar usuários do banco.
    \item executar \texttt{DELETE} sem \texttt{WHERE}.
    \item substituir a linguagem Java.
\end{enumerate}

\item DAO, no contexto do exemplo de Hibernate, aparece como uma classe responsável por:

\begin{enumerate}[label=\alph*)]
    \item persistir o objeto.
    \item criptografar todas as senhas.
    \item apagar o servidor.
    \item substituir o driver JDBC.
\end{enumerate}

\item SQL embutido provê principalmente:

\begin{enumerate}[label=\alph*)]
    \item programação de BD estática.
    \item programação dinâmica via API.
    \item ausência de pré-processamento.
    \item execução apenas em tempo de execução.
\end{enumerate}

\item APIs de acesso a BD oferecem como vantagem:

\begin{enumerate}[label=\alph*)]
    \item não precisar de pré-processamento, sendo mais flexíveis.
    \item checagem completa do SQL em tempo de compilação.
    \item impossibilidade de conexão.
    \item ausência de funções.
\end{enumerate}

\item Uma desvantagem de APIs de acesso a BD é que:

\begin{enumerate}[label=\alph*)]
    \item a checagem do SQL é feita em tempo de execução.
    \item nunca permitem SQL dinâmico.
    \item exigem sempre pré-processador.
    \item não podem usar strings.
\end{enumerate}

\item SQL/CLI significa:

\begin{enumerate}[label=\alph*)]
    \item Call Level Interface.
    \item Command Line Internet.
    \item Create Logic Index.
    \item Control Language Instance.
\end{enumerate}

\item SQL/CLI faz parte:

\begin{enumerate}[label=\alph*)]
    \item do padrão SQL.
    \item apenas do HTML.
    \item do sistema operacional.
    \item de uma linguagem sem relação com BD.
\end{enumerate}

\item Em SQL/CLI, comandos SQL são:

\begin{enumerate}[label=\alph*)]
    \item criados dinamicamente e passados como parâmetros string.
    \item proibidos dentro de programas.
    \item sempre previamente compilados.
    \item substituídos por imagens.
\end{enumerate}

\item Um dos componentes de SQL/CLI é:

\begin{enumerate}[label=\alph*)]
    \item registro de ambiente.
    \item registro de senha MD5.
    \item registro de arquivo PDF.
    \item registro de página HTML.
\end{enumerate}

\item O registro de conexão em SQL/CLI mantém:

\begin{enumerate}[label=\alph*)]
    \item informações de controle de conexão de um BD em particular.
    \item informações sobre uma senha em texto puro.
    \item informações sobre imagens do sistema.
    \item apenas comentários do programa.
\end{enumerate}

\item Procedimentos armazenados são:

\begin{enumerate}[label=\alph*)]
    \item funções ou procedimentos armazenados localmente e executados pelo SGBD.
    \item comandos que só rodam no navegador.
    \item consultas digitadas apenas manualmente.
    \item chaves estrangeiras.
\end{enumerate}

\item Uma vantagem dos procedimentos armazenados é:

\begin{enumerate}[label=\alph*)]
    \item evitar duplicação quando o mesmo procedimento é usado por várias aplicações.
    \item tornar impossível o reuso por várias aplicações.
    \item impedir execução pelo servidor.
    \item eliminar completamente SQL.
\end{enumerate}

\item Outra vantagem dos procedimentos armazenados é:

\begin{enumerate}[label=\alph*)]
    \item reduzir custos de comunicação ao executar no servidor.
    \item aumentar obrigatoriamente o tráfego entre cliente e servidor.
    \item impedir qualquer chamada externa.
    \item remover o SGBD.
\end{enumerate}

\item Uma desvantagem dos procedimentos armazenados é:

\begin{enumerate}[label=\alph*)]
    \item problemas de portabilidade, pois cada SGBD pode ter sua própria sintaxe.
    \item impossibilidade de serem chamados por aplicações.
    \item ausência de execução no servidor.
    \item impossibilidade de uso com SQL.
\end{enumerate}

\item O comando usado para chamar procedimento ou função armazenada é:

\begin{enumerate}[label=\alph*)]
    \item \texttt{CALL}.
    \item \texttt{FETCH}.
    \item \texttt{DROP}.
    \item \texttt{CONNECT}.
\end{enumerate}

\item SQL/PSM é:

\begin{enumerate}[label=\alph*)]
    \item parte do padrão SQL para escrever módulos armazenados persistentes.
    \item uma API exclusiva de Java.
    \item uma técnica de ataque por strings maliciosas.
    \item uma forma de remover cursores.
\end{enumerate}

\item SQL/PSM adiciona ao SQL:

\begin{enumerate}[label=\alph*)]
    \item construções de programação como ramificações e loops.
    \item apenas comandos \texttt{DROP}.
    \item apenas comandos HTML.
    \item apenas criptografia MD5.
\end{enumerate}

\item SQL Injection é:

\begin{enumerate}[label=\alph*)]
    \item uma ameaça à segurança e integridade de sistemas que interagem com BD via SQL.
    \item uma técnica segura de criptografia de senha.
    \item uma forma normal de banco de dados.
    \item uma biblioteca de Java.
\end{enumerate}

\item SQL Injection geralmente envolve:

\begin{enumerate}[label=\alph*)]
    \item uso de strings maliciosas.
    \item uso obrigatório de chaves primárias.
    \item criação de índices.
    \item normalização até BCNF.
\end{enumerate}

\item No exemplo da aula, a entrada \texttt{111111100 OR TRUE} em uma concatenação SQL é perigosa porque:

\begin{enumerate}[label=\alph*)]
    \item pode transformar a condição em verdadeira para muitas tuplas.
    \item cria automaticamente uma chave estrangeira.
    \item impede qualquer exclusão.
    \item criptografa a senha corretamente.
\end{enumerate}

\item Uma prática que ajuda a reduzir risco de SQL Injection é:

\begin{enumerate}[label=\alph*)]
    \item usar comandos parametrizados, como \texttt{PreparedStatement}.
    \item concatenar diretamente toda entrada do usuário.
    \item aceitar qualquer string sem validação.
    \item usar \texttt{DELETE} sem \texttt{WHERE}.
\end{enumerate}

\item Em dados criptografados no exemplo da aula, o \textit{salt} serve para:

\begin{enumerate}[label=\alph*)]
    \item compor o valor usado no hash da senha.
    \item substituir o nome do usuário.
    \item criar uma chave estrangeira.
    \item abrir conexão JDBC.
\end{enumerate}

\item No exemplo de senha da aula, o valor armazenado no banco é:

\begin{enumerate}[label=\alph*)]
    \item o hash da senha combinada com salt.
    \item a senha em texto puro.
    \item apenas o CPF do usuário.
    \item o comando SQL completo.
\end{enumerate}

\item No login com senha criptografada, o servidor compara:

\begin{enumerate}[label=\alph*)]
    \item o hash calculado da senha informada com o hash armazenado.
    \item a senha em texto puro com o CPF.
    \item a consulta SQL com o nome da tabela.
    \item a chave estrangeira com a chave primária.
\end{enumerate}

\end{enumerate}

\section{Questões de Verdadeiro ou Falso}

\begin{enumerate}[resume,label=\textbf{Q\arabic*.}]

\item Julgue as afirmações:

\begin{enumerate}[label=\alph*)]
    \item ( \quad ) A maioria das interações práticas com BD ocorre por programas de aplicação.
    \item ( \quad ) Interfaces web podem acessar bancos de dados.
    \item ( \quad ) A interface interativa é a única forma prática de acessar um BD.
    \item ( \quad ) Programação de BD é um assunto amplo e com várias técnicas.
\end{enumerate}

\item Julgue as afirmações:

\begin{enumerate}[label=\alph*)]
    \item ( \quad ) Java pode ser uma linguagem hospedeira.
    \item ( \quad ) SQL pode ser sublinguagem de dados.
    \item ( \quad ) SQL embutido mistura comandos SQL com linguagem hospedeira.
    \item ( \quad ) Linguagem hospedeira e SQL são sempre a mesma coisa.
\end{enumerate}

\item Julgue as afirmações:

\begin{enumerate}[label=\alph*)]
    \item ( \quad ) O descompasso de impedância envolve incompatibilidades entre modelos.
    \item ( \quad ) Resultados SQL com várias tuplas podem exigir cursores ou iteradores.
    \item ( \quad ) Tipos de dados podem contribuir para o descompasso de impedância.
    \item ( \quad ) O descompasso de impedância só ocorre em bancos NoSQL.
\end{enumerate}

\item Julgue as afirmações:

\begin{enumerate}[label=\alph*)]
    \item ( \quad ) SQL embutido pode usar \texttt{EXEC SQL}.
    \item ( \quad ) Variáveis compartilhadas podem ser prefixadas com dois pontos.
    \item ( \quad ) Um pré-processador pode separar SQL do código hospedeiro.
    \item ( \quad ) SQL embutido nunca permite consultas.
\end{enumerate}

\item Julgue as afirmações:

\begin{enumerate}[label=\alph*)]
    \item ( \quad ) \texttt{CONNECT} pode abrir conexão com servidor de BD.
    \item ( \quad ) \texttt{SET CONNECTION} pode mudar a conexão ativa.
    \item ( \quad ) \texttt{DISCONNECT} pode fechar uma conexão.
    \item ( \quad ) Apenas uma conexão pode existir em todo o programa.
\end{enumerate}

\item Julgue as afirmações:

\begin{enumerate}[label=\alph*)]
    \item ( \quad ) Cursor é útil para processar múltiplas tuplas.
    \item ( \quad ) \texttt{FETCH} move o cursor para a próxima tupla.
    \item ( \quad ) \texttt{CLOSE CURSOR} indica fim do processamento.
    \item ( \quad ) Cursor é usado apenas para apagar tabelas.
\end{enumerate}

\item Julgue as afirmações:

\begin{enumerate}[label=\alph*)]
    \item ( \quad ) SQL dinâmica permite montar comandos SQL em tempo de execução.
    \item ( \quad ) SQL dinâmica pode receber comandos digitados pelo usuário.
    \item ( \quad ) Consulta dinâmica pode ser complexa.
    \item ( \quad ) SQL dinâmica nunca usa strings.
\end{enumerate}

\item Julgue as afirmações:

\begin{enumerate}[label=\alph*)]
    \item ( \quad ) SQLJ é um padrão para embutir SQL em Java.
    \item ( \quad ) SQLJ é traduzido para Java usando JDBC.
    \item ( \quad ) SQLJ usa iteradores para múltiplas tuplas.
    \item ( \quad ) SQLJ é uma técnica exclusiva de criptografia.
\end{enumerate}

\item Julgue as afirmações:

\begin{enumerate}[label=\alph*)]
    \item ( \quad ) JDBC significa Java Database Connectivity.
    \item ( \quad ) JDBC usa drivers para acessar SGBDs.
    \item ( \quad ) \texttt{ResultSet} representa resultado de consulta.
    \item ( \quad ) \texttt{executeQuery} é usado principalmente para \texttt{INSERT}.
\end{enumerate}

\item Julgue as afirmações:

\begin{enumerate}[label=\alph*)]
    \item ( \quad ) Hibernate facilita o mapeamento objeto-relacional.
    \item ( \quad ) Hibernate pode gerar chamadas SQL.
    \item ( \quad ) Hibernate pode melhorar portabilidade entre SGBDs.
    \item ( \quad ) Hibernate elimina completamente a existência de tabelas.
\end{enumerate}

\item Julgue as afirmações:

\begin{enumerate}[label=\alph*)]
    \item ( \quad ) SQL/CLI é Call Level Interface.
    \item ( \quad ) SQL/CLI faz parte do padrão SQL.
    \item ( \quad ) SQL/CLI usa comandos SQL criados dinamicamente e passados como strings.
    \item ( \quad ) SQL/CLI não possui componentes de controle.
\end{enumerate}

\item Julgue as afirmações:

\begin{enumerate}[label=\alph*)]
    \item ( \quad ) Procedimentos armazenados são executados pelo SGBD.
    \item ( \quad ) Procedimentos armazenados podem reduzir duplicação.
    \item ( \quad ) Procedimentos armazenados podem reduzir custo de comunicação.
    \item ( \quad ) Procedimentos armazenados não possuem problemas de portabilidade.
\end{enumerate}

\item Julgue as afirmações:

\begin{enumerate}[label=\alph*)]
    \item ( \quad ) SQL/PSM permite escrever módulos armazenados persistentes.
    \item ( \quad ) SQL/PSM adiciona ramificações e loops ao SQL.
    \item ( \quad ) SQL/PSM aumenta o poder de SQL.
    \item ( \quad ) SQL/PSM é o mesmo que SQL Injection.
\end{enumerate}

\item Julgue as afirmações:

\begin{enumerate}[label=\alph*)]
    \item ( \quad ) SQL Injection ameaça segurança e integridade.
    \item ( \quad ) SQL Injection pode usar strings maliciosas.
    \item ( \quad ) Concatenar entrada do usuário diretamente em SQL pode ser perigoso.
    \item ( \quad ) PreparedStatement aumenta o risco de SQL Injection.
\end{enumerate}

\item Julgue as afirmações:

\begin{enumerate}[label=\alph*)]
    \item ( \quad ) Senhas podem ser armazenadas como hash com salt.
    \item ( \quad ) No login, pode-se comparar o hash calculado com o hash armazenado.
    \item ( \quad ) Armazenar senha em texto puro é sempre a melhor prática.
    \item ( \quad ) HTTPS aparece no fluxo de envio de dados no exemplo da aula.
\end{enumerate}

\end{enumerate}

\section{Questões de Aplicação Objetiva}

\begin{enumerate}[resume,label=\textbf{Q\arabic*.}]

\item Considere o trecho:

\[
\texttt{String sql = "DELETE FROM EMPLOYEE WHERE ssn=" + ssn;}
\]

Qual é o principal problema desse código?

\begin{enumerate}[label=\alph*)]
    \item Risco de SQL Injection por concatenação direta da entrada.
    \item Uso obrigatório de chave estrangeira.
    \item Ausência de \texttt{GROUP BY}.
    \item Uso de \texttt{SELECT} em vez de \texttt{DELETE}.
\end{enumerate}

\item Considere a entrada:

\[
\texttt{111111100 OR TRUE}
\]

em um comando montado por concatenação. O efeito provável é:

\begin{enumerate}[label=\alph*)]
    \item tornar a condição verdadeira para muitas tuplas.
    \item impedir execução do SQL para sempre.
    \item criar uma nova tabela chamada \texttt{TRUE}.
    \item transformar o comando em \texttt{SELECT}.
\end{enumerate}

\item Qual alternativa representa uma forma mais segura de executar comandos com parâmetros em JDBC?

\begin{enumerate}[label=\alph*)]
    \item Usar \texttt{PreparedStatement} com \texttt{?} e ligação de parâmetros.
    \item Concatenar diretamente qualquer entrada do usuário.
    \item Usar \texttt{DELETE} sem \texttt{WHERE}.
    \item Criar uma string SQL com \texttt{OR TRUE}.
\end{enumerate}

\item Em JDBC, para uma consulta que retorna linhas, qual sequência faz mais sentido?

\begin{enumerate}[label=\alph*)]
    \item criar conexão, criar comando, executar \texttt{executeQuery}, processar \texttt{ResultSet}.
    \item criar \texttt{ResultSet}, executar \texttt{DROP}, abrir conexão no final.
    \item fechar conexão, processar resultado, executar consulta.
    \item executar \texttt{DELETE}, usar \texttt{GROUP BY}, criar driver.
\end{enumerate}

\item Em SQL embutido, qual comando é usado para avançar sobre tuplas de uma consulta com várias linhas?

\begin{enumerate}[label=\alph*)]
    \item \texttt{FETCH}.
    \item \texttt{DROP}.
    \item \texttt{CREATE FUNCTION}.
    \item \texttt{CALL}.
\end{enumerate}

\item Em SQLJ, se o iterador lista tipos e nomes dos atributos, ele é:

\begin{enumerate}[label=\alph*)]
    \item designado.
    \item posicional.
    \item dinâmico.
    \item criptográfico.
\end{enumerate}

\item Em SQLJ, se o iterador lista somente os tipos dos atributos, ele é:

\begin{enumerate}[label=\alph*)]
    \item posicional.
    \item designado.
    \item armazenado.
    \item injetado.
\end{enumerate}

\item Qual técnica é mais diretamente associada ao mapeamento entre classes Java e tabelas relacionais?

\begin{enumerate}[label=\alph*)]
    \item Hibernate.
    \item SQL Injection.
    \item \texttt{FETCH}.
    \item SQLCODE.
\end{enumerate}

\item Qual alternativa descreve melhor um procedimento armazenado?

\begin{enumerate}[label=\alph*)]
    \item Módulo armazenado e executado pelo SGBD.
    \item Consulta digitada apenas no terminal.
    \item Tabela virtual derivada de consulta.
    \item Senha criptografada com MD5.
\end{enumerate}

\item Em um sistema de login com senha armazenada como hash com salt, o banco deve armazenar preferencialmente:

\begin{enumerate}[label=\alph*)]
    \item o hash calculado a partir da senha e do salt.
    \item a senha original em texto puro.
    \item o comando SQL usado no login.
    \item o nome da tabela de usuários.
\end{enumerate}

\item Considere o trecho JDBC:
\[
\texttt{Connection con = DriverManager.getConnection(url, user, pass);}
\]
O objeto \texttt{Connection} representa:

\begin{enumerate}[label=\alph*)]
    \item uma conexão ativa com o banco de dados.
    \item uma linha retornada por consulta.
    \item uma tabela temporária.
    \item uma chave estrangeira.
\end{enumerate}

\item Considere o trecho:
\[
\texttt{PreparedStatement ps = con.prepareStatement("SELECT * FROM Aluno WHERE matricula = ?");}
\]
O símbolo \texttt{?} representa:

\begin{enumerate}[label=\alph*)]
    \item um parâmetro a ser ligado antes da execução.
    \item um erro obrigatório de sintaxe.
    \item uma coluna chamada interrogação.
    \item um comando para apagar a tabela.
\end{enumerate}

\item No trecho:
\[
\texttt{ps.setInt(1, matricula);}
\]
o número \texttt{1} indica:

\begin{enumerate}[label=\alph*)]
    \item a posição do primeiro parâmetro \texttt{?} no comando preparado.
    \item o número de linhas da tabela.
    \item o tipo da chave primária.
    \item o número de conexões abertas.
\end{enumerate}

\item Considere:
\[
\texttt{ResultSet rs = ps.executeQuery();}
\]
Esse método é adequado principalmente para:

\begin{enumerate}[label=\alph*)]
    \item comandos \texttt{SELECT} que retornam linhas.
    \item comandos \texttt{UPDATE} que não retornam linhas.
    \item fechar a conexão.
    \item criar uma classe Java.
\end{enumerate}

\item Para percorrer as linhas retornadas em um \texttt{ResultSet}, usa-se normalmente:

\begin{enumerate}[label=\alph*)]
    \item \texttt{while (rs.next())}.
    \item \texttt{while (con.close())}.
    \item \texttt{DROP RESULTSET}.
    \item \texttt{GROUP BY rs}.
\end{enumerate}

\item Para executar um \texttt{UPDATE}, \texttt{INSERT} ou \texttt{DELETE} em JDBC, o método mais apropriado é:

\begin{enumerate}[label=\alph*)]
    \item \texttt{executeUpdate()}.
    \item \texttt{executeQuery()} obrigatoriamente.
    \item \texttt{next()}.
    \item \texttt{getString()}.
\end{enumerate}

\item Considere o código:
\[
\texttt{String sql = "SELECT * FROM Usuario WHERE login='" + login + "' AND senha='" + senha + "'";}
\]
O principal risco é:

\begin{enumerate}[label=\alph*)]
    \item SQL Injection por concatenação de entrada externa.
    \item violação automática da 1FN.
    \item ausência de \texttt{ORDER BY}.
    \item uso de transação serial.
\end{enumerate}

\item Uma entrada como:
\[
\texttt{' OR '1'='1}
\]
em um login vulnerável pode:

\begin{enumerate}[label=\alph*)]
    \item tornar a condição sempre verdadeira e burlar a autenticação.
    \item criar uma chave primária automaticamente.
    \item impedir qualquer consulta \texttt{SELECT}.
    \item normalizar a tabela de usuários.
\end{enumerate}

\item Qual alternativa reduz melhor o risco de SQL Injection em JDBC?

\begin{enumerate}[label=\alph*)]
    \item Usar \texttt{PreparedStatement} e parâmetros, evitando concatenação direta.
    \item Remover todos os comandos \texttt{WHERE}.
    \item Usar \texttt{Statement} com concatenação livre.
    \item Armazenar senha em texto puro.
\end{enumerate}

\item Considere uma aplicação que precisa garantir que duas atualizações ocorram juntas. Em JDBC, uma abordagem adequada é:

\begin{enumerate}[label=\alph*)]
    \item desativar \texttt{autoCommit}, executar as operações e chamar \texttt{commit()} ou \texttt{rollback()}.
    \item usar sempre \texttt{SELECT DISTINCT}.
    \item fechar a conexão antes das atualizações.
    \item executar cada atualização em bancos diferentes sem transação.
\end{enumerate}

\item No JDBC, se uma exceção ocorre no meio de uma transação manual, a ação adequada para desfazer as alterações parciais é:

\begin{enumerate}[label=\alph*)]
    \item chamar \texttt{rollback()}.
    \item chamar \texttt{getString()}.
    \item chamar \texttt{next()}.
    \item chamar \texttt{GROUP BY}.
\end{enumerate}

\item Em SQL embutido, o uso de cursor é necessário principalmente quando:

\begin{enumerate}[label=\alph*)]
    \item uma consulta retorna múltiplas tuplas e o programa precisa processá-las uma a uma.
    \item o comando é \texttt{DROP TABLE}.
    \item não existe linguagem hospedeira.
    \item a consulta retorna sempre uma constante.
\end{enumerate}

\item Considere a sequência de cursor:
\[
\texttt{DECLARE cursor},\quad \texttt{OPEN},\quad \texttt{FETCH},\quad \texttt{CLOSE}
\]
O comando \texttt{FETCH} serve para:

\begin{enumerate}[label=\alph*)]
    \item obter a próxima tupla do resultado.
    \item apagar o cursor fisicamente do disco.
    \item abrir uma nova conexão JDBC.
    \item executar \texttt{COMMIT} automaticamente.
\end{enumerate}

\item Uma stored procedure é útil quando se deseja:

\begin{enumerate}[label=\alph*)]
    \item armazenar lógica de processamento no SGBD para ser chamada por aplicações.
    \item eliminar todas as tabelas.
    \item substituir qualquer índice por imagem.
    \item impedir transações.
\end{enumerate}

\item Em SQL/PSM, procedimentos e funções armazenadas permitem:

\begin{enumerate}[label=\alph*)]
    \item usar estruturas procedurais próximas ao banco, como variáveis e controle de fluxo.
    \item apenas criar arquivos PDF.
    \item remover a necessidade de SQL.
    \item impedir chamadas por aplicações.
\end{enumerate}

\item Hibernate é associado principalmente a:

\begin{enumerate}[label=\alph*)]
    \item mapeamento objeto-relacional entre classes e tabelas.
    \item detecção de deadlock por grafo de espera.
    \item normalização para BCNF.
    \item ordenação por \texttt{ORDER BY}.
\end{enumerate}

\item O problema conhecido como descompasso de impedância aparece porque:

\begin{enumerate}[label=\alph*)]
    \item linguagens de programação e bancos relacionais organizam dados de formas diferentes.
    \item SQL e Java possuem exatamente o mesmo modelo de dados.
    \item todo banco relacional dispensa tipos.
    \item cursores eliminam objetos.
\end{enumerate}

\item Para armazenar senhas com mais segurança, uma prática melhor que texto puro é:

\begin{enumerate}[label=\alph*)]
    \item armazenar hash da senha com salt apropriado.
    \item armazenar a senha original sem proteção.
    \item colocar a senha no nome da tabela.
    \item usar \texttt{SELECT *} no login.
\end{enumerate}

\item Se uma aplicação compara diretamente a senha digitada com uma senha em texto puro armazenada no banco, o problema principal é:

\begin{enumerate}[label=\alph*)]
    \item exposição grave caso o banco seja vazado.
    \item excesso de normalização.
    \item uso obrigatório de cursor.
    \item conflito \textit{write-write}.
\end{enumerate}

\item Em uma consulta JDBC, chamar \texttt{rs.getString("nome")} normalmente significa:

\begin{enumerate}[label=\alph*)]
    \item recuperar o valor da coluna \texttt{nome} da linha atual do \texttt{ResultSet}.
    \item criar uma coluna chamada \texttt{nome}.
    \item abrir uma conexão com senha \texttt{nome}.
    \item executar \texttt{DELETE} na tabela \texttt{nome}.
\end{enumerate}


\end{enumerate}

\newpage

\section*{Gabarito}

\begin{multicols}{2}
\begin{enumerate}[label=\textbf{Q\arabic*.}]
    \item a
    \item a
    \item a
    \item a
    \item a
    \item a
    \item a
    \item a
    \item a
    \item a
    \item a
    \item a
    \item a
    \item a
    \item a
    \item a
    \item a
    \item a
    \item a
    \item a
    \item a
    \item a
    \item a
    \item a
    \item a
    \item a
    \item a
    \item a
    \item a
    \item a
    \item a
    \item a
    \item a
    \item a
    \item a
    \item a
    \item a
    \item a
    \item a
    \item a
    \item a
    \item a
    \item a
    \item a
    \item a
    \item a
    \item a
    \item a
    \item a
    \item a
    \item a
    \item a
    \item a
    \item a
    \item a
    \item a
    \item a
    \item a
    \item a
    \item a
    \item a
    \item a
    \item a
    \item a
    \item a
    \item a
    \item V, V, F, V
    \item V, V, V, F
    \item V, V, V, F
    \item V, V, V, F
    \item V, V, V, F
    \item V, V, V, F
    \item V, V, V, F
    \item V, V, V, F
    \item V, V, V, F
    \item V, V, V, F
    \item V, V, V, F
    \item V, V, V, F
    \item V, V, V, F
    \item V, V, V, F
    \item V, V, F, V
    \item a
    \item a
    \item a
    \item a
    \item a
    \item a
    \item a
    \item a
    \item a
    \item a
    \item a
    \item a
    \item a
    \item a
    \item a
    \item a
    \item a
    \item a
    \item a
    \item a
    \item a
    \item a
    \item a
    \item a
    \item a
    \item a
    \item a
    \item a
    \item a
    \item a
\end{enumerate}
\end{multicols}

\end{document}
